\documentclass{book}
\usepackage{hyperref}
\usepackage[margin=1in]{geometry}
\usepackage{ulem}
\usepackage{tipa} % Add this line to support \textipa command
\hypersetup{
    colorlinks=true,
    linkcolor=black,
    filecolor=magenta,      
    urlcolor=cyan,
    }

\urlstyle{same}

\title{French Conjugation}
\author{Junzhang Luo}

\begin{document}

\tableofcontents

\section{Day 1}
    De lundi à Vendredi, je vais au travail en voiture, ça prends environ quarante minuits.
    Je suis un programmeur sur l'équipe développement logiciel. 
    Je retourne chez moi à 17 heures et je me douche avant je diné. 
    Après diner, je consacre un à deux heures étudier français.
    En géneral, avant je me couche, je joue du piano ou je joue aux jeu de vidéo avec mes amis.
    Sur le week-end, je me reveille un peu tôt pour prends le petit déjeuner.
    Je vais joue au badminton ou jeu de société. Mais le plus souvent, je apprends quelque chose.

    \subsection{Correction}
        De lundi à vendredi, je vais au travail en voiture, et cela prend environ quarante minutes.
        Je suis programmeur dans l'équipe de développement logiciel.
        Je retourne chez moi à 17 heures et je me douche avant de dîner.
        Après le dîner, je consacre une à deux heures à étuidier le français.
        En général, avant de me coucher, je joue du piano ou je joue à des jeux vidéo avec mes amis.
        \\
        Le week-end, je me réveille un peu tôt pour prendre le petit déjeuner.
        Je vais jouer au badminton ou à des jeux de société. Mais le plus souvent, j'apprends quelque chose.

    \subsection{Remarks}
        \begin{itemize}
            \item Pour prends $\to$ pour prendre (use infinitive after ``pour'')
            \item Après le dîner
            \item Avant de + infinitive
            \item Do not use ``sur'' for weekend, we can say le week-end or pendant le week-end
            \item Il m'arrive de + infinitif
            \item Il m'arrive (souvent / parfois / rarement) de + infinitif
        \end{itemize}

    \subsection{C1}
        Du lundi au vendredi, je me rends au travail en voiture, un trajet qui dure environ quarante minutes.
        Je travaille comme programmeur au sein d'une équipe de développement logiciel. 
        En fin de journée, je rentre chez moi vers 17 heures et je prends une douche avant de dîner. 
        Après le repas, je consacre habituellement une à deux heures à l'étude du français.
        Avant de me coucher, il m'arrive de jouer du piano ou de passer un peu de temps à jouer à des jeux vidéo avec mes amis.
        Le week-end, je me lève un peu plus tôt pour prendre le petit déjeuner, puis je fais du badminton ou des jeux de société.
        Cependant, la plupart du temps, j'essaie surtout d'apprendre quelque chose de nouveau.

\section{Day 2}
        J'étudie à l'université de Waterloo, le campus est grand et très agréable. Il y a une jardin avec plusieurs des rochers dans le centre du campus.
        Il y a deux lacs sur le campus, un à Nord du campus, nommée lac Columnbia, à côte de la centre sportif, nommée Columbia Ice Field.
        L'autre lac est à West du campus, entre le collége United et le collége St.Jeromes. 
        Il y a toujours les oies du Canada ses passent du temps dispersés autour du campus.
        Je suis une partie de la département de mathématiques. 
        Il y a un total de trois l'immeubles sous la département de mathématiques: MC, M3, DC.
        M4 est une nouvelle immeuble en cours de construction pendnat deux années, ce n'est pas encore terminé.
        En générale, je vais au PAC, l'autres centre de sport, pour loisirs de ouvrir.
        Je joue au badminton et au basketball. Je vais rarement à la salle de sport.

    \subsection{Correction}
        J'étudie à l'Université de Waterloo. Le campus est grand et très agréable.
        Il y a un jardin avec plusieurs rochers au centre du campus.
        Il y a aussi deux lacs sur le campus : l'un se trouve au nord du campus et s'appelle le lac Columbia, a côté du centre sportif appelé le Columbi Ice Field.
        L'autre est situé à l'ouest du campus, entre le collège United et le collège St. Jerome's.
        \\
        Il y a toujours des oies du Canada qui passent du temps un peu partout sur le campus.
        Je fais partie du département de mathématiques. 
        Il y a trois bâtiments dans ce département: MC, M3 et DC. M4 est un nouveau bâtiment en construction depuis deux ans et il n'est pas encore terminé.
        En général, je vais au PAC, un autre centre sportif, pour faire des activités de loisir.
        J'y joue au badminton et au basket-ball, mais je vais rarement à la salle de sport.

    \subsection{Remark}
        \begin{itemize}
            \item fixed expression: un peu partout (in many different places)
        \end{itemize}

    \subsection{C1}
        J'étudie à l'Université de Waterloo, dont le campus est à la fois vaste et agréable.
        Au centre, on trouve un jardin parsemé de plusieurs rochers, ce qui lui donne un cadre assez paisible.
        Le campus compte également deux lacs : le lac Columnbia, situé au nord près du Columbia Ice Field, et un autre à l'ouest, entre le collège United et le collège St. Jerome's.
        On y voit souvent des oies du Canada se promener librement. 
        Je fais partie du département de mathématique et, pendant mon temps libre, je vais parfois au PAC pour jouer au badminton ou au basket-ball.

\section{Day 4}
    J'étuidie l'informatique pour mon programme de université. 
    En revanche, j'adore le musique, faire du sports et goûter des plats différentes.
    J'apprend jouer du piano quand j'avais trois ans. Je joue au saxophone commencé le septième grade.
    Dès les sports, je joue au badminton, au table-tennis et au basket-ball.
    Sur le week-end, je vais au centre le sport avec mes amis pour jouer au badminton.
    J'adore explorer les nouvelles technologies et créer les logiciels pour résoudre des problems quotidiennes.

    \subsection{Correction}
        J'étudie l'informatique dans le cadre de mon programme universitaire.
        En revanche, j'adore la musique, faire du sport et goûter des plats différents.
        J'ai appris à jouer du piano quand j'avais trois ans.
        Je joue du saxophone depuis la septième année.
        Côté sport, je joue au badminton, au tennis de table et au basket-ball.
        Le week-end, je vais au centre sportif avec mes amis pour jouer au badminton.
        J'adore explorer de nouvelles technologies et créeer des logiciels pour résoudre des problèmes du quotidien.

    \subsection{Remark}
        \begin{itemize}
            \item dans le cadre de (as part of, in the context of, within the framework of)
            \item du quotidiens (of daily life / of everyday life)
            \item bien que (although)
            \item vont bien au-delà de (go well beyond)
            \item être attiré par quelque chose (to be drawn to)
        \end{itemize}

    \subsection{C1}
        Bien que j'étudie l'informatique à l'université, mes centres d'intérêt vont bien au-delà du cadre académique.
        Je suis particulièrement attiré par la musique, le sport et la découverte culinaire. 
        J'ai appris à jouer du piano dès l'âge de trois ans, et je pratique le saxophone depuis la septième année.
        Sur le plan sportif, je joue notamment au badminton, au tennis de table et au basket-ball, que je pratique souvent avec mes amis le week-end au centre sportif.
        Par ailleurs, j'aime explorer les nouvelles technologies et concevoir des logiciels capables de répondre à des besoins concrets du quotidien.

\section{Day 5}
    À mon avis, la meilleure ville j'avais jamais visiter est Xi'an, la captiale de Shaan Xi province.
    Xi'an est une ville avec la culture et les plats locaux très délicieux. 
    Dans la ville, on se trouve le vielle mur de la ville, le segment plus tôt à dès la dynastie Qin.
    On se monte sur le mur et fait un tour complet de la ville.
    Il y a aussi de tour du tambour et de tour de la cloche.
    Pour la nourriture, il y a plusieurs plats différentes à Xi'an. 
    Parmi des plats célébres, les plus renommés comprenenté l'agneau avec la brioche vapeur, le hamburger chinois et le Bao avec la soupe à l'intérieur.
    
    \subsection{Correction}
        À mon avis, la meilleure ville que j'aie jamais visitée est Xi'an, la capitale de la province du Shaanxi.
        Xi'an est une ville riche en culture et en spécialités locales très délicieuses.
        Dans cette ville, on trouve les anciens remparts, dont les premières parties remontent à la dynastie Qin.
        On peut monter sur les remparts et faire le tour complet de la ville.
        Il y a aussi la tour du Tambour et la tour de la Cloche.
        Côté nourriture, il existe de nombreux plats différents à Xi'an.
        Parmi les plats célèbres, les plus connus comprennent la soupe de mouton avec pain émietté, le hamburger chinois et les bao farcis avec de la soupe à l'intérieur.

    \subsection{Remarks}
        \begin{itemize}
            \item remonter à (go back to)
            \item after parmi, use definite group
            \item Côté + noun (as for / when it comes to / on the ... side / regarding)
        \end{itemize}

    \subsection{C1}
        À mes yeux, Xi'an est l'une des ville les plus remarquables que j'aie jamais visitées.
        Capitale de la province du Shaanxi, elle se distingue non seulement par la richesse de son patrimoine historique, mais aussi par la diversité de sa culture locale.
        La ville est notamment connue pour ses impressionnants remparts anciens, sur lesquels on peut monter afin d'admirer la ville tout en en faisant le tour.
        Elle abrite également des monuments emblématiques, tels que la tour du Tambour et la tour de la Cloche.
        Enfin, Xi'an est réputée pour sa gastronomie, avec de nombreuses spécialités locales particulièrement savoureuses.

    \subsection{Remarks}
        \begin{itemize}
            \item l'une des \dots que j'aie jamais + participe passé
            \item riche en + noun
            \item on trouve (describe places)
            \item les remparts (for wall)
            \item être répouté pour / être connu pour
            \item se distinguer par + noun (to stand out because of)
            \item non seulement ... mais aussi ... (not only but also)
            \item patrimoine (heritage)
            \item monter sur quelque chose, so we use sur + lesquels
            \item tout en + present participle (while + verb-ing)
            \item faire le tour de quelque chose (to go around something / go all the way around something)
            \item the second en is a pronoun replacing des remparts
            \item abriter (to house, to contain)
            \item emblématiques (iconic, emblematic, symbolic)
            \item tel / telle / tels / telles que (such as)
        \end{itemize}

\section{Day 6}
    Est-il important d'apprendre plusieurs langues ?
    À mon avis, apprendre de plusieurs langues servi seulment les avantages donc il est important l'apprendre.
    D'abord, savez comme parler une langue fournit des opportunités parler aux gens avec cela contexte linguistique. 
    On peut faire le gen qui parle cette langue sentir la chaleur et faire nouveau ami.
    Ensuite, on apprend leur contexte de la culture et comprend plus choses de les quotidenne.
    Cependant, il prend beaucoup de temps et l'effort à apprendre une langue et atteindre un niveau où on peut parler à quelqu'un qui parler cette langue maternelle.
    Mais pendant le processus de apprendre une nouvelle langue, on non seulement apprend comme la parler, mais aussi obtenit plusieurs des connaissances des pays où la langue est appliquée.
    Dernièrement, savez plus langues on permettre à voyager à les pays différents et essayer, et voir plusieurs chose.
    En conclusion, apprendre plusieurs langues est important et avantageuses.

    \subsection{Correction}
        À mon avis, apprendre plusieurs langues présente seulement des avantages, donc il est important de le faire.
        \\
        D'abord, savoir parler une langue offre des occasions de parler avec des personnes issues de ce contexte linguistique.
        On peut faire en sorte que les gens qui parlent cette langue se sentent accueillis et se faire de nouveaux amis.
        \\
        Ensuite, on découvre leur culture et on comprend davantage leur vie quotidienne.
        \\
        Cependant, apprendre une langue demande beaucoup de temps et d'efforts afin d'atteindre un niveau où l'on peut parler avec quelqu'un dont c'est la langue maternelle.
        \\
        Mais, pendant le processus d'apprentissage, on apprend non seulement à parler la langue, mais aussi beaucoup de choses sur les pays où elle est utilisée.
        \\
        Enfin, connaître plusieurs langues permet de voyager dans différents pays et d'y vivre des expériences plus riches.
        \\
        En conclusion, apprendre plusieurs langues est important et avantageux.

    \subsection{C1}
        À mon avis, l'apprentissage de plusieurs langues constitue une richesse incontestable, tant sur le plan humain que culturel.
        \\
        Tout d'abord, la maîtrise d'une langue permet d'entrer plus facilement en contact avec des personnes d'un autre horizon linguistique,
        tout en créant une relation plus chaleureuse et plus authentique.
        Elle favorise ainsi les échanges et facilite la création de nouveaux liens d'amitié.
        \\
        En outre, apprendre une langue permet de mieux comprendre la culture, les habitudes et le quotidien de ceux qui la parlent.
        \\
        Certes, cet apprentissage exige du temps, de la persévérance et des efforts considérables avant d'atteindre un niveau suffisant pour converser aisément avec un locuteur natif.
        \\
        Cependant, ce parcours est particulièrement enrichissant, car il permet non seulement d'acquérir une nouvelle langue, mais aussi d'approfondir sa connaissance des pays où elle est employée.
        \\
        Enfin, le plurilinguisme facilite les voyages et permet de vivre des expériences plus riches et plus authentiques.
        \\
        En conclusion, apprendre plusieurs langues représente un atout majeur dans le monde d'aujourd'hui.
        
    \subsection{Remarks}
        \begin{itemize}
            \item tant sur le plan X que Y (both on the X level and the Y level)
            \item entrer en contact avec (to get in touch with)
            \item horizon (literally means horizon, but can be background, world, sphere)
            \item en outre (moreover)
            \item ceux (those)
            \item certes (admittedly, certainly)
            \item exiger (to require)
            \item aisément (easily)
            \item approfondir sa connaissance de (to deepen one's knowledge of)
            \item un atout (an asset)
            \item représenter un atout (to be an important advantage, to represent a real asset)
        \end{itemize}
    




\end{document}